%% bab2.tex
%% BAB 2 - KAJIAN PUSTAKA
%% Format sesuai Panduan Penyusunan Tesis ITS 2021

\chapter{KAJIAN PUSTAKA}

Pada bagian ini, teori, temuan, dan bahan penelitian sebelumnya yang diperoleh dari berbagai referensi yang dijadikan dasar melakukan penelitian yang diusulkan ini dibahas. Hal yang relevan dengan subjek/topik/\textit{state of the art} yang diteliti saja yang diuraikan.

\section{Kajian Pustaka}

Kajian pustaka merupakan rangkuman singkat yang komprehensif tentang semua materi terkait yang terdapat di dalam berbagai referensi. Bagian ini bisa merupakan tampilan diskusi atau debat antar pustaka. Selain itu juga bisa menjelaskan tentang teknik/kaidah/peralatan atau teknologi yang akan dan/atau telah digunakan dalam melaksanakan penelitian.

Uraian yang ditulis diarahkan untuk menyusun kerangka pendekatan atau konsep yang diterapkan dalam penelitian. Materi yang disampaikan diusahakan dari referensi terbaru dan asli, misalkan dari \textit{journal papers}. Perlu dicatat bahwa minimal 30\% dari total pustaka di dalam kajian pustaka adalah berasal dari artikel jurnal ilmiah yang relevan.

\subsection{Topik Kajian Pertama}

Uraikan topik kajian pustaka pertama yang relevan dengan penelitian. Sertakan kutipan dari referensi yang digunakan menggunakan sistem Harvard \citep{contoh2024}.

\subsection{Topik Kajian Kedua}

Uraikan topik kajian pustaka kedua yang relevan dengan penelitian.

\subsection{Topik Kajian Ketiga}

Uraikan topik kajian pustaka ketiga yang relevan dengan penelitian.

\section{Dasar Teori}

Dasar teori merupakan semua teori yang diambil/dipilih berdasarkan kajian pustaka yang melatarbelakangi permasalahan penelitian/studi yang akan/sedang dilakukan. Dasar teori juga akan digunakan sebagai pedoman untuk mengerjakan penelitian lebih lanjut. Bentuk dasar teori bisa berupa uraian kualitatif atau model/persamaan matematis.

\subsection{Teori Pertama}

Jelaskan teori pertama yang menjadi dasar penelitian ini.

\subsection{Teori Kedua}

Jelaskan teori kedua yang menjadi dasar penelitian ini. Jika diperlukan, cantumkan rumus/persamaan matematis. Contoh penulisan rumus adalah sebagai berikut:

\begin{equation}
    E = mc^2
    \label{eq:einstein}
\end{equation}

\noindent dengan:
\begin{tabular}{@{}lcl@{}}
    $E$ & = & energi (Joule) \\
    $m$ & = & massa (kg) \\
    $c$ & = & kecepatan cahaya ($3 \times 10^8$ m/s)
\end{tabular}

\vspace{0.5cm}

Penulisan rumus dimulai pada batas kiri, rata dengan kalimat tepat di atasnya. Nomor rumus disesuaikan dengan nomor bab letak rumus tersebut dan ditulis di ujung kanan baris tersebut (\textit{aligned right}).

\subsection{Teori Ketiga}

Jelaskan teori ketiga yang menjadi dasar penelitian ini.

\section{Penelitian Terkait}

Bagian ini berisi review dari penelitian-penelitian terkait yang telah dilakukan sebelumnya. Jelaskan kelebihan dan kekurangan dari masing-masing penelitian serta posisi penelitian ini dibandingkan dengan penelitian sebelumnya.

Tabel~\ref{tab:penelitian-terkait} menunjukkan perbandingan penelitian terkait.

\begin{table}[htbp]
    \caption{Perbandingan Penelitian Terkait}
    \label{tab:penelitian-terkait}
    \centering
    \begin{tabular}{|l|l|l|l|}
        \hline
        \textbf{Peneliti} & \textbf{Metode} & \textbf{Kelebihan} & \textbf{Kekurangan} \\
        \hline
        Peneliti A (2023) & Metode X & Kelebihan A & Kekurangan A \\
        \hline
        Peneliti B (2022) & Metode Y & Kelebihan B & Kekurangan B \\
        \hline
        Peneliti C (2021) & Metode Z & Kelebihan C & Kekurangan C \\
        \hline
    \end{tabular}
    
    \vspace{0.3cm}
    \footnotesize{Sumber: Hasil Kajian Pustaka}
\end{table}
